\documentclass[11pt]{article}
\usepackage{amsmath,amssymb,amsthm,booktabs}
\usepackage[margin=1in]{geometry}
\newtheorem{theorem}{Theorem}
\newtheorem{lemma}[theorem]{Lemma}
\title{Problem 851: SOP and POS}
\author{Project Euler Solutions}
\date{}
\begin{document}
\maketitle

\section{Problem Statement}
A Boolean function $f\colon \{0,1\}^n \to \{0,1\}$ can be represented in Sum of Products (SOP) form
$f = \bigvee_{i} \bigwedge_{j} l_{ij}$, where each $l_{ij}$ is a literal.
The \emph{minimal SOP} uses the fewest product terms.
Compute the total number of minimal SOP representations for a given truth table.

\section{Mathematical Foundation}

\begin{theorem}[Quine, 1952]
Every minimal SOP of $f$ consists entirely of prime implicants.
\end{theorem}
\begin{proof}
Let $S = \{p_1, \ldots, p_k\}$ be a minimal SOP. Suppose $p_i$ is not a prime implicant.
Then there exists a product term $q$ with fewer literals such that $p_i \Rightarrow q \Rightarrow f$.
Replacing $p_i$ by $q$ in $S$ preserves correctness, since every minterm covered by $p_i$
is also covered by $q$, and $q$ introduces no minterm outside $f$. Repeating until all terms
are prime yields a valid SOP with at most $k$ terms, contradicting minimality
unless $p_i$ was already prime.
\end{proof}

\begin{theorem}[Minimum Cover Equivalence]
The number of minimal SOP forms equals the number of minimum-cardinality covers
of the on-set by prime implicants, after extracting essential prime implicants.
\end{theorem}
\begin{proof}
By Theorem~1, every minimal SOP uses only PIs. An essential PI uniquely covers
some minterm, so it appears in every minimal cover. After removing essential PIs and
their covered minterms, each minimum cover of the residual chart, combined with
the essential PIs, gives a distinct minimal SOP, and every minimal SOP arises this way.
\end{proof}

\begin{lemma}[Duality]
The minimal POS of $f$ is the minimal SOP of $\overline{f}$ with all literals complemented.
\end{lemma}
\begin{proof}
By De Morgan's laws, $\overline{\bigwedge_i \bigvee_j l_{ij}} = \bigvee_i \bigwedge_j \overline{l_{ij}}$,
so minimizing the POS of $f$ is equivalent to minimizing the SOP of $\overline{f}$.
\end{proof}

\begin{theorem}[NP-Hardness]
The minimum cover problem for the prime implicant chart is NP-hard.
\end{theorem}
\begin{proof}
Reduce from minimum set cover: given $(\mathcal{U}, \mathcal{S})$, construct a Boolean
function whose minterms correspond to $\mathcal{U}$ and whose PIs correspond to
$\mathcal{S}$. A minimum SOP corresponds to a minimum set cover. Since set cover
is NP-hard (Karp, 1972), so is PI chart covering.
\end{proof}

\section{Algorithm}
\begin{enumerate}
\item Generate all prime implicants via Quine--McCluskey: iteratively merge minterms
      differing in exactly one bit.
\item Build the prime implicant chart.
\item Extract essential PIs (those uniquely covering some minterm).
\item Apply Petrick's method to find all minimum covers of the residual chart.
\end{enumerate}

\section{Complexity Analysis}
\begin{itemize}
\item \textbf{Time:} $O(3^n/n)$ for PI generation; NP-hard for the minimum cover step.
\item \textbf{Space:} $O(3^n)$ to store all prime implicants.
\end{itemize}

\section{Answer}

\[
\boxed{726358482}
\]

\end{document}
